\documentclass[a4paper,11pt]{article}

\usepackage[utf8]{inputenc}
\usepackage[T1]{fontenc}
\usepackage[spanish]{babel}
\usepackage{geometry}
\usepackage{hyperref}
\usepackage{xurl}
\usepackage{array}

\geometry{margin=2.3cm}
\hypersetup{
    colorlinks=true,
    urlcolor=blue,
    linkcolor=black
}

\title{Guia de archivos de la entrega}
\author{Juan Jose Prado Neira \\ Javier Rego Bargos}
\date{}

\begin{document}

\maketitle

\section*{Contenido de la entrega}

Esta carpeta contiene los archivos principales del TFM.

\begin{itemize}
    \item \texttt{circuito\_segmentation/}: parte usada para preparar el dataset,
    entrenar el modelo de segmentacion y guardar modelos/resultados.
    \item \texttt{socket\_vision/}: programas usados para conectar el coche con el
    portatil, procesar vision en tiempo real y mandar comandos a la Jetson.
    Dentro contiene otro \texttt{README.md} para navegar esa parte.
    \item \texttt{paper/}: articulo academico de dos paginas.
    \item \texttt{presentacion/}: presentacion y guion de exposicion.
\end{itemize}

\section*{Archivos entregados por OneDrive}

Por problemas de espacio, las carpetas mas pesadas se entregan en una carpeta
compartida de OneDrive.

\begin{center}
\url{https://nubeusc-my.sharepoint.com/:f:/g/personal/juanjose_prado_rai_usc_es/IgC97ixKN1sLQa5HmM_D72O4AWIi2pscvEqPPGk-5HM1bmw?e=KOsI5p}
\end{center}

Dentro de esa carpeta de OneDrive estan:

\begin{itemize}
    \item \path{circuito_segmentation/data/}: dataset dividido en entrenamiento
    y validacion.
    \item \path{circuito_segmentation/outputs/}: predicciones, mascaras y
    resultados visuales generados durante las pruebas.
    \item \path{circuito_segmentation/checkpoints/}: modelos entrenados en
    formato PyTorch, incluyendo el mejor checkpoint usado para exportar a ONNX.
    \item \path{videos/}: videos recopilados durante el desarrollo y las
    pruebas del sistema.
\end{itemize}

\section*{Notas}

El flujo de entrenamiento esta documentado en:

\begin{center}
\texttt{circuito\_segmentation/notebooks/entrenamiento\_completo\_makesense.ipynb}
\end{center}

El flujo principal del coche se lanza desde:

\begin{center}
\texttt{socket\_vision/conexion/scripts/main.sh}
\end{center}

\end{document}
